\chapter{Results and Analysis}
\label{ch:results}

This section reports the verification evidence for the framework requirements.
The results of the two case studies represent the deliverable of the fifth objective of this project.

\section{User tests}
\label{section:user-trial-results}

The user tests assess the time required to maintain requirements, retrieve supporting evidence, assess changes and save mission inputs through the framework.
The protocol in Table~\ref{tab:framework-mop-protocols} assigns ten tasks to each participant, with a limit of 6~min per task and a total allocation of 60~min.
Each task uses the same prepared initial state for both participants.
The mission conditions and numerical allocations are purely illustrative.
They provide controlled starting points for measuring the effort required to perform only the assigned engineering activities, not an entire mission definition.

The tests use the GUI and a prepared mission environment.
For FR-09, participants are also given access to a terminal.
Before that task, they are instructed to use \texttt{git status --short} to view repository status and \texttt{git diff} to inspect changes, and practice both commands on the prepared baseline.
Preparation and this briefing occur before timing begins.
Task time runs from presentation of the task to submission.
Submission is when the participant declares verbally that they have completed the task, after which a screenshot is taken.
Only the screenshots from the student club member are presented in this document, since the deterministic analyses of the framework produced the same outputs under equal inputs for both participants.

Each MOP requires submissions to match the expected results from the fixed inputs and compares its two elapsed times with the limit separately, without averaging.
Table~\ref{tab:user-trial-first-task} shows the first task, and Appendix~\ref{appendix:user-trials} gives the prepared inputs, expected results and complete table of user test results.

% Keep the evidence layout shared so the first row and its appendix copy agree.
\newlength{\UserTrialBodyWidth}
\newcommand{\UserTrialImageMaxHeight}{265pt}
\newcommand{\UserTrialImage}[1]{%
\vspace{0pt}\includegraphics[width=\linewidth,height=\UserTrialImageMaxHeight,keepaspectratio]{#1}}
\newenvironment{UserTrialEvidence}{%
\footnotesize
\setlength{\tabcolsep}{3pt}%
\setlength{\parskip}{0pt}%
\renewcommand{\arraystretch}{1.08}%
\setlength{\UserTrialBodyWidth}{\dimexpr\linewidth-10\tabcolsep\relax}%
\begin{tabular}{@{}>{\raggedright\arraybackslash}p{0.14\UserTrialBodyWidth}>{\centering\arraybackslash}p{0.30\UserTrialBodyWidth}>{\centering\arraybackslash}p{0.30\UserTrialBodyWidth}>{\centering\arraybackslash}p{0.09\UserTrialBodyWidth}>{\centering\arraybackslash}p{0.09\UserTrialBodyWidth}>{\centering\arraybackslash}p{0.08\UserTrialBodyWidth}@{}}
\toprule
\textbf{MOP/Task} & \textbf{Shared initial state} & \textbf{Final state of student club member} & \textbf{\shortstack[t]{Time\\taken by\\PhD\\student}} & \textbf{\shortstack[t]{Time\\taken by\\student\\club\\member}} & \textbf{Verifies FR?} \\
\midrule
}{\bottomrule\end{tabular}}
\newcommand{\UserTrialRow}[5]{%
\textbf{MOP-#1}\newline FR-#1\newline #2 &
\UserTrialImage{fig/user-trials/mop#1-initial.png} &
\UserTrialImage{fig/user-trials/mop#1-club.png} & #3 & #4 & #5 \\[5pt]}
\newcommand{\UserTrialFirstRow}{\UserTrialRow{01}{Create and save a requirement}{226.88}{243.76}{Yes}}

\begin{sidewaystable}[p]
\centering
\caption[First user test task.]{First user test task. Screenshots show the shared initial state and the student club member's submission. Times are in seconds.}
\label{tab:user-trial-first-task}
\begin{UserTrialEvidence}
\UserTrialFirstRow
\end{UserTrialEvidence}
\end{sidewaystable}


All 20 submissions in the reported final test set matched the expected results, so no task failed this criterion.
Both times for each of the ten reported user test MOPs are below the limit.
Their sums are 32~min 37.44~s for the PhD role and 34~min 14.84~s for the club role.
Individual times range from 88.20 to 318.16~s.
The longest time, for MOP-06 in the club role, is 41.84~s below the per-task limit.

The 60~min allocation reflects the duration of the stakeholders' weekly meeting.
Verification of all ten user test MOPs by both stakeholders would suggest that they can assess the impact of requirement changes and review updated EPS results within that meeting.
This is considerably faster than what was observed during stakeholder analysis, when stakeholders would have to wait until the next weekly meeting to see the results of tasks comparable to those considered in the user tests.

Furthermore, the fact that a single engineering tool was used across requirements, analyses and component comparisons mitigates the core problem of evidence fragmentation.
The framework stores the relationships between requirements, design inputs and supporting results in one model, which reduces the effort needed to reconstruct those relationships during a review and supports evaluating alternatives on a common basis.
The Git commands run in the terminal provide the configuration control required for FR-09, including inspection of changes and retrieval of an authorized baseline.

The user tests were repeated with the same participants, and the reported results come from the last set of tests.
Earlier attempts were halted because results did not match the expected outputs.
Participants also found the GUI so unclear that explanations of its use consumed the entire allocated time.
These attempts provided qualitative feedback for GUI improvements.
Repeated practice and GUI improvements can explain why the final set had no such failures.
Repetition is considered valid because the framework is expected to have a learning curve.
The reported times therefore characterize performance after practice.
A discussion of the GUI's evolution and detailed user feedback is outside the scope of this document.

\section{Catalogue inspection}
Table~\ref{tab:cots-catalog-counts} reports MOP-10 by the five generic CSRM EPS hardware types.
The catalogue contains 147 distinct products, each with a source datasheet link.
Combined products count once under each type they inherit.
All five hardware types meet the minimum of five products, so MOP-10 is verified.

\begin{table}[htbp]
\centering
\caption{MOP-10 catalogue inspection by inherited CSRM hardware type. Counts overlap for combined products.}
\label{tab:cots-catalog-counts}
\small
\begin{tabular}{lrl}
\toprule
CSRM hardware type & Distinct products & Minimum of five met? \\
\midrule
Battery & 15 & Yes \\
SolarArray & 12 & Yes \\
RegulatorAndConverter & 28 & Yes \\
PowerCircuit & 78 & Yes \\
ProcessorAndMemory & 19 & Yes \\
\midrule
Distinct products across all types & 147 & \\
Products with a source datasheet link & 147 & \\
\bottomrule
\end{tabular}
\end{table}

\section{EPS case studies}
\label{section:eps-case-study-results}

The EPS case studies verify the calculations required by FR-12 and FR-14 to FR-20 using BIRDS-4 and OreSat data.
Mission publications supply component ratings and selected calculation results \citep{birds2021batteryverification,shrestha2022birds4powerbudget,greenberg2021oresat,oresat2025solardesign}.
Explicit assumptions supply the remaining inputs.

\subsection{Verification setup}

The BIRDS electrical case uses a two-cell 3G30A solar section and a 3S2P nickel-metal hydride (NiMH) battery.
The OreSat case uses a two-cell XTE-SF section and two independent 2S1P lithium-ion (Li-ion) packs.
Both baseline sections face the Sun.
The comparisons also cover changes in illumination, cell technology, series and parallel arrangements, bus voltage and power-transfer method.
Appendix~\ref{appendix:eps-case-studies} groups the inputs and selected results by the requirement they verify.

Independent reference values were calculated directly from the input data, without calling the framework's numerical functions.
The reference summed component masses and volumes, calculated duty-weighted state powers, and obtained phase energies from mean power and duration.
Loss-adjusted phase energies determined source-power demand.
Battery sizing used eclipse demand, the assumed cycle-life/DoD relation and the energy margin.
Power and energy ratios were rounded upward to obtain integer cell counts.
These sums, products and ratios used decimal arithmetic, with 40 significant digits for products and ratios.
A separate orbital reference solved Kepler's equation at uniformly spaced times and tested Earth shadow in Cartesian coordinates.
This provided an independent check of the framework's shadow-boundary and panel-projection calculations.

Each nonzero numerical output had to agree with its reference within 1\%.
Zero values, integer counts and discrete outcomes required exact agreement.
The appendix reports maximum deviations and representative comparisons, with repeated outputs and deliberately failing cases omitted from the detailed tables.

\subsection{Results and discussion}

The numerical comparisons met the 1\% criterion for every assessed requirement in both mission cases.
Integer counts and discrete outcomes also matched their references.
No unexpected mismatch was found.
Appendix~\ref{appendix:eps-case-studies} separates the evidence for mass accounting, sizing, component choices, electrical configurations, conversion, photovoltaic characteristics, transfer paths and protection.

The largest deviation was 0.331\% in the BIRDS photovoltaic power comparison for FR-18.
The framework calculated $5.123\times0.050=0.25615$~W, compared with the published 0.257~W label.
The next largest deviation was 0.315\% for FR-14, where the OreSat cell ratings gave 2.34738~W against a published two-cell value of 2.34~W.
Both comparisons use printed values with limited precision.
They check consistency with the publications at the stated tolerance, while the separately calculated references provide a more precise check of the arithmetic.

The largest deviation from an independent calculation was 0.000864\% in the BIRDS battery-energy requirement.
The framework returned 4243.934~J against 4243.970~J from the reference.
This difference follows from the finite time grid used to estimate eclipse duration in the orbital reference.
Doubling that grid changed eclipse duration by at most 0.00205\%.
The minimum integer cell counts remained unchanged.
The small energy difference therefore had no effect on the sizing decisions checked here.

These results support using the framework for EPS accounting, preliminary sizing and comparisons of entered component options.
The evidence includes independent calculations of both the component quantities and the orbital geometry that links illumination to the power budget.
For FR-15 and FR-16, the technology and arrangement comparisons reproduced the expected changes in power, voltage, energy and capacity.
For FR-17, FR-19 and FR-20, terminal currents, transfer-path powers and protection comparisons agreed with the separately evaluated electrical relations.
This connects the numerical agreement to the decisions the framework is intended to support.

The remaining uncertainty concerns specific input models.
The complete verification curves use three photovoltaic points, and the battery life fits are assumed.
The published BIRDS 1.575~W measured peak remains unreconstructed.
OreSat's partial solar sweep supplies power-product checks but lacks the endpoints needed for a complete characteristic.
The case studies consequently provide stronger evidence for correct calculation and comparison than for reproducing either mission's measured energy performance.


\section{Economic Analysis}
\label{section:economic-analysis}

The analysis follows Ulrich and Eppinger's sequence of a base case, sensitivity analysis, development tradeoffs and qualitative assessment \citep[Chap.~17]{ulrich2013diseno}.
The comparison is between using the framework and continuing the as-is workflow with two funded researchers.
ASCEND does not sell services, so the monetary benefit considered here is the value assigned to personnel time released for other research.
Cash expenditure is assessed separately.
The modeled period comprises four months of development followed immediately by three mission years, following the CUBICS mission lifespan.
All amounts are held at constant 2026 Canadian dollar values.

\subsection{Base case and included costs}

ASCEND is assumed to provide a scholarship of CA\$1500 per month for four months devoted entirely to framework development.
The resulting scholarship expense is CA\$6000.
An OpenAI subscription at CA\$27.68 per month supported GUI development and COTS parameter extraction from datasheets during the same four months.
This initial subscription expense totals CA\$110.72.
The combined direct development cash expenditure is:
\begin{equation}
    C_{\mathrm{direct}}=4(1500)+4(27.68)=6110.72.
    \label{eq:economic-direct-cost}
\end{equation}
Of this CA\$6110.72, CA\$6000 is assigned to ASCEND and CA\$110.72 was paid from personal funds.

Initial supervision and user test participation are covered by existing funded positions.
Initial supervision is excluded from the effort valuation.
The two researchers and the personnel responsible for subsequent maintenance are each valued at CA\$1500 per month for a 35~h working week, comprising seven hours from Monday to Friday.
Using 52 weeks and 12 months per year gives the hourly funding proxy:
\begin{equation}
    c_h=\frac{12(1500)}{52(35)}\approx9.89\ \text{CA\$/h}.
    \label{eq:economic-hourly-value}
\end{equation}
This proxy values research capacity at its stipend funding cost.
Three user tests failed because participants found the GUI unclear.
These tests consumed 3~h of combined participant time.
One participant also had two tests halted after 11 and 23~min because results did not match expectations.
The other participant had one test halted after 3~min for the same reason.
The final successful tests are assigned two combined hours to include instructions and preparation.
Total test participation is 5~h 37~min, with an economic cost of:
\begin{equation}
    C_{\mathrm{test}}=c_h\left(3+\frac{11+23+3}{60}+2\right)\approx55.55.
    \label{eq:economic-test-cost}
\end{equation}
This funded effort is included as an opportunity cost at the end of development month four.
Total economic development costs are CA\$6166.27, with no additional cash payments for test participation.
The base case assumes that all released time is used for other productive research.
Stipend payments remain unchanged in both workflows.

\subsection{Review effort and maintenance}

The effort baseline uses a comparable update delivered two days later by two participants who reported working only on that activity.
The base case assigns two full working days of seven hours to each person and treats that instance as typical.
This gives 28 combined hours per review cycle.
Section~\ref{section:user-trial-results} gives times of 32~min 37.44~s and 34~min 14.84~s for the ten reported tasks.
For this estimate, each participant's reported task time is rounded up to one hour to include instructions and preparation.
A comparable cycle with the framework therefore takes two combined hours and releases 26 hours of personnel time.

Comparable work occurred weekly during the four-month observation period.
The forecast assumes the same frequency for the first 18 mission months, covering Phases 0 and A.
One cycle per month is assumed for the next six months of Phase B.
No further review cycles are credited in the third mission year.
Weekly activity is distributed across months using $52/12$ cycles per month, with no separate reduction for holidays or breaks.
These assumptions give 78 weekly cycles and six monthly cycles, releasing 2184 hours before maintenance.

Maintenance and database expansion consume ten combined personnel hours per month during the first nine mission months.
This effort assumes AI assistance for COTS parameter extraction.
One subscription at CA\$27.68 per month is therefore included during those nine months.
This initial period covers the first year's maintenance.
One further month of ten combined hours and one subscription payment is assigned to the end of mission year two.
Mission year three has no maintenance or subscription costs.
Total support effort is 100 hours, and the ten additional subscription payments total CA\$276.80.
Net released time during mission operations is 2084 hours.

\subsection{Net present value}

A 7\% annual real discount rate is adopted from the Treasury Board of Canada Secretariat's guidance for Canadian regulatory cost-benefit analysis \citep[Sec.~6.1]{tbs2019costbenefit}.
It serves as a Canadian reference rate for this assessment.
The equivalent monthly rate is $i=(1.07)^{1/12}-1$.
The valuation date is the start of the internship, and all expenditure and capacity benefits are placed at month end.
Development comprises four cash payments of CA\$1527.68 and the test participation cost $C_{\mathrm{test}}$ valued at the end of month four.
For mission month $m$, let $H_m$ denote the released hours before maintenance, $M_m$ the maintenance hours and $A_m$ the additional subscription expense.
The economic net present value is:
\begin{equation}
    \begin{aligned}
        V_{\mathrm{econ}}={}&-\sum_{t=1}^{4}\frac{1527.68}{(1+i)^t}
        -\frac{C_{\mathrm{test}}}{(1+i)^4}\\
        &+\sum_{m=1}^{36}\frac{c_h(H_m-M_m)-A_m}{(1+i)^{m+4}}.
    \end{aligned}
    \label{eq:economic-npv}
\end{equation}
Table~\ref{tab:economic-npv} gives the period totals using the unrounded hourly proxy and discount rate.
The base case gives an economic net present value of approximately CA\$12661.

\begin{table}[htbp]
\centering
\small
\caption{Economic value by period in constant 2026 Canadian dollars. Mission years start after the four-month internship. Present values use monthly discounting, and totals use unrounded values.}
\label{tab:economic-npv}
\begin{tabular}{@{}lrrrr@{}}
\toprule
\textbf{Period} & \textbf{\shortstack{Released\\time (h)}} & \textbf{\shortstack{Maintenance\\time (h)}} & \textbf{\shortstack{Net economic\\value (CA\$)}} & \textbf{\shortstack{Present\\value (CA\$)}} \\
\midrule
Development & 0 & 0 & $-6166.27$ & $-6079.62$ \\
Mission year 1 & 1352 & 90 & 12232.20 & 11522.21 \\
Mission year 2 & 832 & 10 & 8101.99 & 7218.11 \\
Mission year 3 & 0 & 0 & 0.00 & 0.00 \\
\midrule
Total & 2184 & 100 & 14167.92 & 12660.69 \\
\bottomrule
\end{tabular}
\end{table}

The cash assessment credits no income or payroll reduction from the released time.
The continuing stipends are common to both workflows and cancel from the incremental cash comparison.
Development and subsequent subscription payments total CA\$6387.52, with a cash net present value of $-$CA\$6285.77.
The positive economic result represents the assigned value of additional research capacity.

\subsection{Sensitivity and interpretation}

The baseline effort is recalled from one instance, and the task times were measured after earlier unsuccessful user tests with prepared inputs.
There is no controlled comparison of work effort under the two workflows.
The estimate therefore depends on the assumed transfer of that effort difference to repeated mission reviews.
Holding other assumptions fixed, the economic net present value reaches zero at approximately 5.74 baseline hours per person and cycle, compared with the assumed 14 hours.
A separate sensitivity calculation uses the same 14-hour baseline and varies the fraction of released time credited at the hourly proxy.
The economic result reaches zero when approximately 36.5\% of that capacity receives its full assigned value, while all development and support costs remain included.
Using the alternative 3\% social rate discussed in the Treasury Board guide gives approximately CA\$13491 \citep[Sec.~6.1]{tbs2019costbenefit}.

The zero yearly cost for C3 in Table~\ref{tab:framework-candidate-scores} is interpreted as zero software license expenditure.
The AI subscription and personnel effort are separate support costs.
No avoided subscription payments for C1, C2 or C4 are credited, since the comparison uses the as-is workflow.
This valuation includes development expenditure, which is a sunk cost for subsequent support decisions \citep[Chap.~17, App.~A]{ulrich2013diseno}.

The EPS case studies in Section~\ref{section:eps-case-study-results} demonstrate detection of infeasible configurations before hardware fabrication.
No avoided hardware purchase, fabrication or paid service is monetized, and no residual value is credited at the end of the assessment.
The additional value of reusable models, traceable evidence and student experience remains qualitative.
CUBICS 2026 Stream~2 offers a maximum grant of CA\$350000 over up to three years \citep[Sec.~7.1]{csa2026cubics}.
The modeled cash expenditure is approximately 1.83\% of that ceiling.
The positive economic value supports an internal capability investment under the stated effort, utilization and support assumptions.
